% --- Plik: chapter3_design.tex ---

\chapter{Projekt i realizacja systemu}
\label{ch:projekt} % To jest unikalny znacznik tego rozdziału

Niniejszy rozdział stanowi zasadniczą część merytoryczną pracy dyplomowej. Autor szczegółowo opisuje w nim strukturę projektową urządzeń, układów lub algorytmów oraz ich praktyczną realizację. 

Treść powinna w sposób ustrukturyzowany dokumentować proces realizacji projektu od założeń wstępnych do efektu końcowego, uwzględniając:
\begin{itemize}
	\item opracowanie koncepcji oraz uzasadnienie wyboru stosowanych metod, narzędzi, technologii lub elementów składowych,
	
	\item dokumentację projektową i inżynierską (np. schematy, modele matematyczne, diagramy blokowe, algorytmy),
	
	\item przebieg prac implementacyjnych, realizacyjnych lub symulacyjnych oraz integrację poszczególnych modułów i podsystemów.
\end{itemize}

\noindent \textbf{Stylistyka:} tekst należy redagować w bezosobowej formie czasu przeszłego (np. \textit{„zaprojektowano”, „zaimplementowano”, „przeprowadzono dobór”}), zachowując charakter sprawozdania technicznego.

